%%%%

%%%   Самарский государственный технический университет

%%%   Кафедра прикладной математики и информатики

%%%

%%%   Преамбула для статьи в журнал "Вестник Самарского государственного

%%%   технического университета. Серия: «Физико-математические науки»"

%%%   Ориентирована под MikTEx 2.4 for Win32

%%%

%%%   Хотя сам журнал собирается под GNU/Linux на дистрибутиве TeXLive



\documentclass[11pt,a4paper,twoside]{article}



\usepackage[cp1251]{inputenc}

\usepackage[T2A]{fontenc}

\usepackage[english,russian]{babel}

\usepackage{samgtubib}

\usepackage[dvips]{graphicx}

\usepackage[multiple]{footmisc}



\usepackage{samgtu}

\renewcommand{\VestnikYear}{2009} % Год издания Вестника

\renewcommand{\VestnikNumber}{2\,(19)} % Номер Вестника

\renewcommand{\VestnikMonth}{Ноябрь} % Месяц выхода Вестника

\renewcommand{\VestnikData}{01.11.09} % Дата подписания Вестника в печать

\renewcommand{\VestnikUPL}{26,64} % Усл. печ. л.

\renewcommand{\VestnikUIL}{25,73} % Уч.-изд. л.

\renewcommand{\VestnikRegNumber}{479/09} % Регистрационный номер

\renewcommand{\VestnikOrder}{994} % Номер заказа






%\usepackage{nccboxes}
%\usepackage{rotating_pr}
%
%\begin{document}
%
%\newcommand{\totop}[1]{\raisebox{6pt}[1pt][2pt]{#1}} 
%\newcommand{\tobot}[1]{\raisebox{-6pt}[1pt][2pt]{#1}} 

\renewcommand{\firstpage}{125}
\renewcommand{\lastpage}{135}

\udk{519.233.5}
\msc{62G08}


\title{Построение моделей регрессионного типа для~описания \mbox{теплового} состояния системы двух~плоских тел в~режиме циклического контакти\-рования}
{Построение моделей регрессионного типа \ldots}

\author{В.\,В.~Стулин}
{С\,т\,у\,л\,и\,н~В.\,В.}

\institute{\samgtu.}

\email{E-mail}{stulinvv@mail.ru}

\otherinf{\fullauthor{Владимир Васильевич Стулин}\regalia{к.т.н., доц.}\position{доцент}
\dept{каф. высшей математики и прикладной информатики}}

\keywords{циклический контактный теплообмен, система
интегральных уравнений, планирование вычислительного эксперимента, модели
регрессионного типа}

\workabstract{Сложные аналитические решения задач циклического
контактного теплообмена в форме систем интегральных уравнений
приведены к критериальному виду и экономным численным анализом
преобразованы в полиномиальные модели на основе применения
методологии планирования эксперимента. Аппроксимация искомых
функций осуществлялась по дискретным точкам с использованием
формул Бонне.
Вычисления показали достаточно быструю сходимость приближений и
в практических расчётах, как правило, использовалось 7--11 итераций.
Получены 13 критериальных уравнений регрессионного типа,
содержащие наиболее важные и разнородные по составу и структуре
образования характеристики квазиустановившейся стадии циклического
контактного теплообмена. Оценка адекватности моделей выполнена с~помощью множественного коэффициента корреляции.}

\engtitle{Regression Models Construction for Describing the Thermal System State of Two Flat Bodies in Cyclic Contact}

\engauthor{V.\,V.~Stulin}{S\,t\,u\,l\,i\,n~V.\,V.}

\enginstitute{\engsamgtu.}

\otherenginfm{ \fullauthor{Vladimir  V. Stulin} \regalia{Ph.\,D.
(Techn.)} \position{Associate Professor} \dept{Dept. of Higher
Mathematics and Applied Computer Science}}

\engabstract{Sophisticated analytical solutions of cyclic contact heat transfer problems in the form of integral equations were reduced to criteria form and converted into polynomial models based on the application of experiment planning methodology with economical numerical analysis. Approximation of the desired functions was performed by discrete points using Bonnet formulas, calculations showed quite rapid convergence of approximations and in practical calculations the number of iterations was 7--11. Thirteen criteria equations of regression type were received; the equations contain the most important and diverse in composition and formation structure characteristics of quasi-steady stage of cyclic contact heat exchange. The evaluation of the adequacy of models was made using multiple correlation coefficient.}

\engkeywords{cyclic contact heat exchange, system of integral equations, numerical experiment planning, models of  the regression type}

\date{21/V/2012}
\revisiondate{25/VII/2012}

\maketitle


Вопросы управления и оптимизации большой группы процессов
технологической теплофизики (это прежде всего процессы горячей
обработки металлов давлением) связаны с необходимостью решения
задач циклического контактного теплообмена (ЗЦКТ) с чередующимися
во времени краевыми условиями (IV$\rightleftarrows$II) (или
(IV$\rightleftarrows$III), (IV$\rightleftarrows$I)) рода на
контактной поверхности. В задачах подобного класса автоматически
возникают вопросы учёта и описания отдельно или совместно
действующих разнородных тепловых источников как в зонах
контактирования объектов, так и в зонах их пространственной
аккумуляции. К последним следует прежде всего отнести тепловые
потоки $q_v$ от трения тел по контактной поверхности, тепловые
потоки $q_w$ от пластического формоизменения отдельных элементов и
тепловые потоки $q_f$ от начального теплосодержания элементов
контактной системы. В~аналитическом плане некоторые из указанных
выше задач для произвольного цикла контактирования и~квазиустановившейся стадии циклического контактного теплообмена
(ЦКТ) решены в [1--3].

Однако для численной проработки и последующего анализа указанные решения в
том виде, в котором они представлены после интегрирования соответствующих
краевых задач теплопроводности, обладают целым рядом принципиальных
недостатков. В частности, итоговые решения ЗЦКТ получены не в традиционной
функциональной форме, содержащей явную зависимость теплофизических
характеристик от исследуемых факторов, а в виде систем интегральных
уравнений Фредгольма.

Соответственно, возникает необходимость использования основных положений
безразмерного параметрического анализа \cite{stulin:bib4} в рамках
реализуемого вычислительного эксперимента с одновременным
привлечением для решения поставленных задач методов теории
планирования эксперимента [5--9]. Эти методы широко и успешно применяются в области чисто
экспериментальных исследований. Вместе с тем априори можно
полагать, что в качестве числовой информации, которая
закладывается в матрицу планирования, могут использоваться не
только экспериментальные данные, но и результаты пассивного
численного анализа аналитических решений.

С учётом изложенного применим методику активного вычислительного
эксперимента для описания температурного режима контактной системы
<<металлозаготовка~-- пограничный слой~-- инструмент>> при
действии циклического теплового источника $q_f$ от нагретой
металлозаготовки \cite{stulin:bib1}. С этой целью запишем систему
двух интегральных уравнений Фредгольма относительно искомых
функций $f_{s0}^{(\infty )} (x)$ и $q_z^{(\infty )} (t)$,
полученную в \cite{stulin:bib1} для квазиустановившейся стадии ЦКТ
(число циклов контактирования $m\to \infty )$ при действии
теплоисточника $q_f$:
\begin{itemize}


\item температурное распределение в инструменте перед началом
контактирования:
\begin{equation}
\label{stulin:eq1}
f_{s0}^{(\infty )} (x) =\frac{2a_s }{\lambda _s
}\sum\limits_{n=1}^\infty \frac{U_n^2 C_n (x)}{\exp \left(a_s
U_n^2 \tau ^+ \right)-1} \int_0^{\tau _c } q_z^{(\infty )}
(\tau) \exp \left(a_s U_n^2 \tau \right)d\tau,
\end{equation}
где
$$
C_n (x)  =\frac{\lambda _s U_n \cos U_n (l_s -x)+\alpha _s \sin
U_n (l_s -x)}{U_n^2 \left[ \lambda _s l_s U_n \cos U_n l_s +(
\lambda _s +\alpha _s l_s)\sin U_n l_s \right]};
$$


\item тепловой поток в зоне контактирования:
\begin{multline}
\label{stulin:eq2} q_z^{(\infty )} (t)=\frac{2a_s}{\lambda _s }\sum\limits_{n=1}^\infty \mu _n \sin \mu _n \left\{ f_z^0 \left(
\frac{H^\ast }{\mu _n }\cos \frac{\mu _n }{H^\ast }-\frac{1}{{\tt Bi}_s
}\sin \frac{\mu _n }{H^\ast }\right)\right. + \\ 
+ \frac{1}{l_s} \int _0^{l_s } f_{s0}^{(\infty )} (\xi)
 \left[ \frac{\mu _n}{H^\ast {\tt Bi}_s} \cos \frac{\mu _n}{H^\ast }\left( 1-\frac{\xi }{l_s } \right)+ \right. \\ + \left. \left.
 \sin \frac{\mu _n
}{H^\ast }\left(1-\frac{\xi }{l_s } \right) \right]d\xi \right\}
\exp \left(-\frac{\mu _n^2 }{H^{\ast 2}}{\tt Fo}_s  \right).
\end{multline}

\end{itemize}

В уравнениях \eqref{stulin:eq1}, \eqref{stulin:eq2} введены следующие обозначения: $\mu _n$, $U_n $ "--- корни соответствующих
трансцендентных уравнений, полученных в \cite{stulin:bib1}; индекс
$z$ присваивается всем величинам, относящимся к заготовке, индекс
$s$ "--- инструменту; $x$, $t$ "--- пространственная и временная
переменные, ${\tt Fo}_s = a_s t/ l^2_s$ "--- критерий Фурье.

Остальные величины являются постоянными в процессе ЦКТ: $\lambda_s$, $\lambda _z$, $a_s$,  $a_z$, $l_s$, $l_z$ "--- коэффициенты
теплопроводности, температуропроводности и толщина инструмента и
заготовки соответственно; длительность цикла $\tau^+=\tau _c +\tau
_r$ состоит из контактного периода $\tau _c $ и неконтактного
периода (паузы)~$\tau _r $; $\alpha _s $ "--- коэффициент
поверхностного теплообмена между инструментом и окружающей средой
(по плоскости $x=l_s$); $\alpha ^+=1/\rho $ "--- условный
коэффициент теплообмена в зоне контакта ($x=0$) между инструментом
и заготовкой в период контактирования $\tau _c $ (выражается через
термическое сопротивление $\rho$ пограничного слоя между
контактирующими телами); %$K_l =l_z/l_s$, $K_a = a_z/a_s$, 
$H^\ast = (l_z/l_s)\sqrt{a_s/a_z}$ "---
безразмерный критерий, характеризующий темп перестройки
температурной обстановки в инструменте по отношению к заготовке;
${\tt Bi}_s =\alpha _s / (\lambda _s l_s )$, ${\tt Bi}_z^+ = \alpha ^+ / (\lambda _z l_z )$ "--- критерии Био.

Для общности теплофизического анализа выполним его в пространстве
безразмерных комплексов \cite{stulin:bib4}. С этой целью систему
интегральных уравнений (\ref{stulin:eq1}), (\ref{stulin:eq2})
запишем в безразмерном виде:
\begin{multline}
\label{stulin:eq3}
 \hat {f}_{s0}^{(\infty )} (\delta _s K_f )_{\min}^r =2\frac{K_\varepsilon
}{H^\ast }\sum\limits_{n=1}^\infty \frac{{U_n}{{\tt Bi}_s^{-1}}\cos
\left[ U_n \left( 1-\delta _s K_f \right) \right]+\sin \left[U_n
(1-\delta _s K_f ) \right]}{U_n^2 \left[ {U_n }{{\tt Bi}_s ^{-1}}\cos U_n +\left(1+{{\tt Bi}_s ^{-1}} \right)\sin U_n \right]}\times \\
 \times \sum\limits_{K_q =0}^{K_q^\ast -1} {\tt Ki}_z^{(\infty )} \left( \Delta {\tt Fo}_s^c K_q  \right)
 \bigl({1-\exp ( -U_n^2 {\tt Fo}_s^+ )}\bigr)^{-1}
%\times \\
%\times 
\bigl(\exp \{ U_n^2 [ \Delta {\tt Fo}_s^c ( K_q +K_q^b )-{\tt Fo}_s^c ]\}-\\
-\exp [ U_n^2 (\Delta {\tt Fo}_s^c K_q -{\tt Fo}_s^c ) ] \exp (-U_n^2 {\tt Fo}_s^r ) \bigr),
 \end{multline}
\begin{multline}
\label{stulin:eq4} {\tt Ki}_z^{(\infty )} ( \Delta {\tt Fo}_s^c K_q )=
2\sum\limits_{m=1}^\infty \psi_f^{-1} ( \mu _n ) \sin \mu _n 
\Bigl\langle f_z^0 \Bigl( \frac{H^\ast }{\mu _n }\cos \frac{\mu _n }{H^\ast }-\frac{1}{{\tt Bi}_s }\sin \frac{\mu _n }{H^\ast } \Bigr)+  \\
 +2\sin \frac{\mu _n \delta _s K_f^b }{2H^\ast }\sum\limits_{K_f =0}^{K_f^\ast -1} \hat {f}_{s0}^{(\infty )} (\delta _s K_f )_{\min}^r  
 \Bigl\{ \frac{H^\ast }{\mu _n } \sin \frac{\mu _n }{H^\ast }\Bigl[1-\delta _s \Bigl( \frac{K_f^b }{2}+K_f \Bigr) \Bigr]+ \\
 +\frac{1}{{\tt Bi}_s }\cos \frac{\mu _n }{H^\ast } \Bigl[ 1-\delta _s \Bigl( \frac{K_f^b }{2}+K_f \Bigr) \Bigr] \Bigr\} \Bigr\rangle 
 \exp \Bigl( -\mu _n^2 \frac{\Delta {\tt Fo}_s^c K_q }{H^{\ast 2}}\Bigr).
 \end{multline}
В уравнениях \eqref{stulin:eq3}, \eqref{stulin:eq4} введены следующие
обозначения: $U_n$, $\mu _n $ "--- корни соответствующих
характеристических уравнений, полученных в \cite{stulin:bib1},
явные выражения этих уравнений и функции $\psi _f^{-1} (\mu _n )$
ввиду их громоздкости здесь не приводятся; $K_\varepsilon = (\lambda _z/\lambda _s) \sqrt {a_s /a_z } $ "--- безразмерный критерий, характеризующий тепловую активность заготовки по отношению к
инструменту; индексы $f$, $q$ относятся к величинам,
дискретизация которых осуществляется по пространственной
переменной и времени соответственно; 
$\hat {f}_{s0}^{(\infty )} (\delta _s K_f )_{\min }^r =(f_s^{(\infty )} (\delta _s K_f )_{\min }^r )/ f_z^0$ "--- безразмерная минимальная температура по сечению инструмента в момент начала контактирования
на квазиустановившейся стадии; $f_z^0 $ "--- начальная постоянная
температура заготовки; $\delta _s =\Delta l_s =l_s / K_f^\ast$
"--- длина шага разбиения толщины инструмента $l_s $ при общем
числе интервалов разбиения $K_f^\ast =10$ $( K_f =0, 1, 2, \ldots, K_f^\ast)$; 
${\tt Ki}_z^{(\infty )} ( \Delta {\tt Fo}_s^c K_q )=\bigl(q_s^{(\infty )} ( \Delta {\tt Fo}_s^c K_q )l_z \bigr) \times (\lambda _z f_z^0 )^{-1}$ "--- критерий Кирпичёва,
содержащий поток взаимного теплообмена $q_s^{(\infty )} \left(
\Delta {\tt Fo}_s^c K_q \right)$ между заготовкой и инструментом;
$\Delta {\tt Fo}_s^c = (a_s \Delta \tau _c )/ l_s^2$ "--- длина (критериальная) шага разбиения временного этапа
контактирования ${\tt Fo}_s^c$ при общем числе интервалов разбиения $K_q^\ast =10$ $( K_q =0, 1, 2, \ldots, K_q^\ast )$; $\Delta \tau _c =\tau _c/K_q^\ast $; ${\tt Fo}_s^c =a_s \tau _c / l_s^2$, ${\tt Fo}_s^r =a_s \tau _r / l_s^2 $, ${\tt Fo}_s^+ = a_s \tau ^+ / l_s^2 $ "--- безразмерные числа Фурье.

Специально подобранные фиксированные числа $K_f^b$, $K_q^b $
(числа Бонне) согласно применяемым формулам Бонне
\cite{stulin:bib10} для соответствующих десяти основных интервалов
разбиения с шагом $\delta _s$, $\Delta {\tt Fo}_s^c $ устанавливались
численно (в нашем случае "--- из требования удовлетворения
условиям неидеального контакта между заготовкой и инструментом в
центре плана выполняются условия $0<K_f^b <1$, $0 < K_q^b < 1$).

Дополнительно по результатам решения системы уравнений
\eqref{stulin:eq3}, \eqref{stulin:eq4} рассчитывались температурные поля в отдельных элементах контактной системы как в~период контакта, так и во время паузы, которые затем учитывались при составлении итоговых регрессионных моделей. Например, для инструмента "--- наиболее важного по соображениям эксплуатационной
надежности элемента при известном тепловом потоке $q^{(\infty )}(t)$ температурные поля определялись по следующим формулам:
\begin{itemize}
\item[а)] для интервала контактирования 
\begin{multline}
\label{stulin:eq5}
 T_c^\ast (x,t_c )=\frac{2a_s}{\lambda _s}\sum\limits_{n=1}^\infty U_n^2
C_n (x)\exp (-a_s U_n^2 t_c ) \times \\ \times 
\left[  \int_0^{t_c } q^{(\infty )}(\tau ) \right.\exp \left (a_s U_n^2 \tau \right)d\tau - \\
 -\left[ 1-\exp \left(a_s U_n^2 \tau ^+ \right) \right]\left.
^{-1} \int_0^{\tau _c } q^{(\infty )}(\tau )\exp \left(a_s
U_n^2 \tau \right)d\tau  \right], \quad t_c \in [0,\tau _c ];
\end{multline}

\item[б)] для неконтактного периода 
\begin{multline}
\label{stulin:eq6} T_r^\ast (x,t_r )=\frac{2a_s }{\lambda _s
}\sum\limits_{n=1}^\infty U_n^2 C_n (x)\exp \left[ -a_s U_n^2 \left( \tau _c +t_r  \right)\right]
\times \\
 \times 
\left[
1-\exp (-a_s U_n^2 \tau ^+) \right]^{-1} \int_0^{\tau _c } q(\tau )\exp \left(a_s U_n^2 \tau
 \right)d\tau,
\quad t_r \in [\tau _c ,\tau ^+].
\end{multline}

\end{itemize}


Система (\ref{stulin:eq3}), (\ref{stulin:eq4}) решалась методом
последовательных приближений, процесс итераций заканчивался при
максимальном 3\%-ном расхождении двух последовательных приближений
величины критерия ${\tt Ki}_z^{(\infty )} \left(\Delta {\tt Fo}_s^c K_q
\right)$. Аппроксимация искомых функций $\hat {f}_{s0}^{(\infty )}
(\delta _s K_f )_{\min}^r$, ${\tt Ki}_z^{(\infty )} \left( \Delta {\tt Fo}_s^c
K_q  \right)$ осуществлялась по дискретным точкам с использованием
формул Бонне \cite{stulin:bib10}, которые в совокупности
составляют содержание второй теоремы о среднем значении и
эффективно используются для приближенного вычисления интегралов от
произведения функций, принадлежащих классу $C [0,\infty )$.
Вычисления показали достаточно быструю сходимость приближений, и
в практических расчётах, как правило, использовалось 7--11 итераций.

Следует особо подчеркнуть, что в работе с целью поиска более
эффективного и экономного численного алгоритма использовались
различные варианты обобщения формул Бонне. Обобщение основных
формул Бонне \cite{stulin:bib10} на дискретный промежуток $[x_0, x_n ]$ $(x_i<x_{i+1}; i=0, 1, \ldots, n-1)$, содержащий $n$ частей, позволило получить две практически важные
аппроксимационные формулы, которые из всех рассмотренных вариантов
по числу производимых итераций оказались для нашего случая
наиболее приемлемыми. Приведём указанные формулы.

\begin{itemize}
\item[1.] Для монотонно убывающей положительной функции $f(x)\ge 0$ (этому условию соответствует температурное распределение в~инструменте при \mbox{$m\to \infty $}) и интегрируемой $\varphi (x)$ имеем
\begin{multline}
\label{stulin:eq7}
  \int_{x_0 }^{x_n } f(x)\varphi (x)dx=
\sum\limits_{i=0}^{n-1}  \int_{x_i }^{x_{i+1} } f(x)\varphi
(x)dx  = \\ =\sum\limits_{i=0}^{n-1} f(x_i )\left[ \Phi (x_0
+\Delta x(i+\bar {K}))-\Phi (x_i ) \right],\quad  \Phi
'(x)=\varphi (x),
\end{multline}
$\xi_{i+1}=x_0+\Delta x(i+\bar{K})$  $(i=0, 1, \ldots, n-1)$ "---
внутренние точки Бонне \cite{stulin:bib10}; $x_i \le \xi _{i+1} \le x_{i+1}$;
$\Delta x=(x_n -x_0 )/n$, $0\le K_{i+1} \le 1$ "--- числа Бонне для каждого
$(i+1)$-го интервала; $ \mathop{\min}\limits_i K_{i+1} \le \bar
K \le \mathop{\max}\limits_i K_{i+1}$ $(\bar {K}={\rm const})$
"--- число Бонне для всего промежутка $[x_0, x_n ]$.

\item[2.] Для монотонно возрастающей положительной функции $f(x)\ge 0$
(этому условию соответствует возрастающий во времени тепловой
поток в зоне контакта заготовки с инструментом при $m\to \infty$)
и интегрируемой $\varphi (x)$ имеем
\begin{multline}
\label{stulin:eq8}
  \int_{x_0}^{x_n} f(x)\varphi (x)dx=
\sum\limits_{i=0}^{n-1} \int_{x_i }^{x_{i+1}} f(x)\varphi
(x)dx =\\= \sum\limits_{i=1}^n f(x_i)\left[\Phi (x_i )-\Phi (x_0
+\Delta x(i-1+\bar {K})) \right],
\end{multline}
$\xi_i=x_0+\Delta x(i-1+\bar{K})$
$(i=1, 2, \ldots, n)$ "--- внутренние точки Бонне; $x_{i-1} \le \xi _i \le x_i$; $\Delta x=(x_n -x_0
)/n$, $0\le K_i \le 1$ "--- числа Бонне для каждого
$i$-того интервала; $\mathop {\min}\limits_i K_i \le \bar 
K \le \mathop {\max}\limits_i K_i$ $(\bar {K}={\rm const})$ "--- число Бонне для всего промежутка $[x_0, x_n ]$.


\end{itemize}

При расчётах в качестве варьируемых факторов были задействованы все безразмерные комплексы (критерии), входящие в правые части
уравнений \eqref{stulin:eq3}, \eqref{stulin:eq4} и других уравнений в критериальной форме с~параметрами оптимизации $y_1$, $y_2$, $\dots$, $y_{13} $, явное описание которых приводится ниже. 
Измеряемые физические и технологические параметры процесса, содержащиеся в~критериях, а также области их достаточно широкого изменения заимствованы из различных источников [11--13] и представлены в~табл.~1. 
Предельные значения коэффициентов теплоотдачи $\alpha$, указанные в табл.~1, не могут быть достигнуты в условиях
реального технологического процесса, поэтому имеют теоретическое значение: при $\alpha ^+\to \infty $ возникает идеальный тепловой контакт между контактирующими телами, а при $\alpha _s \to \infty $ температура внешней поверхности инструмента ($x=l_s$) совпадает с температурой охлаждающей среды. На основе приведенных в табл.~1
данных вычислялись границы изменения выбранных факторов ЦКТ в~безразмерной форме.

\begin{table}[t!] 

\setlength{\tabcolsep}{4.5pt}
\centering \small
\caption{\bf Пределы варьирования физических и~технологических
параметров}\vspace{-3mm}
\begin{tabular}{c|c|c|c|c|c|c|c|c|c|c} \hline
\footnotesize Параметры & \footnotesize $a_z \cdot 10^{5}$  & \footnotesize $a_s \cdot 10^{5}$  & \footnotesize $\lambda_z$ & \footnotesize $\lambda_s$ & \footnotesize $l_z \cdot 10^{3}$ & \footnotesize $l_s \cdot 10^{3}$  & \footnotesize $\alpha ^+$ & \footnotesize $\alpha _s$ & \footnotesize $\tau_c$  & \footnotesize $\tau ^+ $ \\
\hline
\footnotesize Размерность & \multicolumn{2}{c}{\footnotesize м$^2$/с} & \multicolumn{2}{|c|}{\footnotesize Вт/(м$\cdot $К)}
& \multicolumn{2}{c}{\footnotesize м} & \multicolumn{2}{|c|}{\footnotesize Вт/(м$^2 \cdot $К)}
& \multicolumn{2}{c}{\footnotesize с}  \\
\hline
\rule{0mm}{12pt}%
min & 0,3 & 0,49&16,75&16,75&5&2&4187&4,2&1&5 \\
max & 0,6 & 0,98 & 33,5 & 33,5& 20& 10& $\infty $& $\infty $& 30& 50 \\
\hline
\end{tabular}
\end{table}

Величины, выбранные в качестве независимых переменных (варьируемых факторов
$x_1$, $x_2$, $\ldots$, $x_6 $), а также их уровни (согласно общепринятой схеме планирования экспериментов [5--9]) приведены в табл.~2.

Следует особо отметить, что факторы $x_2 =1/{\tt Bi}_s$, $x_3
=K_\varepsilon$, $x_4 =H^\ast$, $x_6 ={\tt Fo}_s^+ $ входят в систему
уравнений (\ref{stulin:eq3}), (\ref{stulin:eq4}) в явном виде, а
факторы $x_1 =1/{\tt Bi}_z^+ $, $x_5 =\tau $ "--- в неявном виде,
так как $x_1 $ входит одновременно в функцию $\psi _f (\mu _n )$ и
характеристическое уравнение для определения корней $\mu _n $ через
соотношение ${\tt Bi}_s^+ =K_\varepsilon / (H^\ast {\tt Bi}_z^+ )$.
Фактор $x_5 $ содержится в критериях ${\tt Fo}_s^r ={\tt Fo}_s^+ (1-\tau
)$, ${\tt Fo}_s^c =\tau {\tt Fo}_s^+ $, входящих в систему уравнений
(\ref{stulin:eq3}), (\ref{stulin:eq4}).

Численным значениям критериев Био ${\tt Bi}_s $ и ${\tt Bi}_z^+ $ соответствуют знаменатели приведенных в табл.~2 дробных выражений (столбцы 2 и 3 таблицы), в числителе указанных дробных выражений приведены соответствующие критериям
Био значения безразмерных термических сопротивлений $x_1$, $x_2$.
Предельные значения критериев Био (${\tt Bi}_s \to \infty$, ${\tt Bi}_z^+ \to \infty $)
имеют аналогичный теплофизический смысл, что и параметры $\alpha ^+$, $\alpha _s $ в~табл.~1.


\begin{sidewaystable}
%\hspace{1.5cm}
%\begin{minipage}{15.5cm}
%\begin{table}[h]
\small\centering
\caption{\bf Условия проведения численного эксперимента для задачи циклического
контактного теплообмена с учётом инициирования теплоты нагретой
металлозаготовки} \vspace{-3mm}
\begin{tabular}{l|c|c|c|c|c|c} \hline
 & 
\footnotesize $x_1= \left({\tt Bi}_s^+ \right)^{-1}$ & 
\footnotesize $x_2=({\tt Bi}_s)^{-1}$ & 
\footnotesize $x_3=K_\varepsilon$ & 
\footnotesize $x_4=H^*$ &  
\footnotesize $x_5=\tau =\tau_c/\tau^+$& 
\footnotesize $x_6={\tt Fo}_s^+$ \\

\multicolumn{1}{c|}{\footnotesize Факторы}  & 
\footnotesize (термическое & 
\footnotesize (термическое & 
\footnotesize (критерий теп- & 
\footnotesize (температурно- & 
\footnotesize (скважность & 
\footnotesize (критериальная\\
\multicolumn{1}{c|}{\footnotesize (безразмерные} & 
\footnotesize сопротивление & 
\footnotesize сопротивление & 
\footnotesize ловой актив- & 
\footnotesize геометрический & 
\footnotesize процесса цик- & 
\footnotesize длительность \\
\multicolumn{1}{c|}{\footnotesize числа   и кри-}  &
\footnotesize  в зоне контак-        & 
\footnotesize  на внутренней & 
\footnotesize ности контакти- & 
\footnotesize критерий контак-    & 
\footnotesize лического кон- & 
\footnotesize единичного \\
\multicolumn{1}{c|}{\footnotesize терии)}& 
\footnotesize тирования) & 
\footnotesize  поверхности & 
\footnotesize рующих   тел) & 
\footnotesize тирующих тел) &  
\footnotesize тактирования) & 
\footnotesize цикла) \\
& 
\footnotesize    & 
\footnotesize  инструмента)& 
& &  & \\
%              & $x_1= \left({\rm Bi}_s^+ \right)^{-1}$ & $x_2=({\rm Bi}_s)^{-1}$ & тел, $x_3=K_\varepsilon$ & тел, $x_4=H^*$ &  $x_5=\tau =\tau_c/\tau^+$&
%              \\ 
              \hline
\rule{0mm}{12pt}%
Основной  &  \tobot{$\dfrac{0,8}{1,25}$} & \tobot{$\dfrac{2000}{0,0001}$} &      &      &    & \\
уровень, ($x_{i0}$) & & & \totop{2,035} & \totop{9,27} & \totop{0,4}& \totop{61,37} \\
\rule{0mm}{12pt}%
Интервалы    & \tobot{$\dfrac{0,34}{2,94}$} & \tobot{$\dfrac{841}{0,00119}$}  &      &      &    & \\
варьирования, ($\Delta x_i)$ &  &  & \totop{0,67} & \totop{3,71} & \totop{0,08} & \totop{25,71} \\
\rule{0mm}{12pt}%
Верхний   & \tobot{$\dfrac{1,14}{0,88}$} & \tobot{$\dfrac{2841}{0,00035}$} &      &      &    & \\
уровень, $(+1)$ &  & & \totop{2,705}& \totop{12,98}& \totop{0,48} & \totop{87,08}  \\
\rule{0mm}{12pt}%
Нижний    & \tobot{$\dfrac{0,46}{2,17}$} & \tobot{$\dfrac{1159}{0,00086}$} &      &      &    & \\
уровень, ($-1$) &  & & \totop{1,365}& \totop{5,56} & \totop{0,32}& \totop{35,86}  \\
\rule{0mm}{12pt}%
Верхняя звёздная  & \tobot{$\dfrac{1,6}{0,625}$}  & \tobot{$\dfrac{4000}{0,00025}$} &      &      &    & \\
точка, $(+2,378)$ &  & & \totop{3,62}& \totop{18,08} & \totop{0,6}& \totop{122,5}  \\
\rule{0mm}{12pt}%
Нижняя  звёздная  & \tobot{$\dfrac{0,0}{\infty}$}  & \tobot{$\dfrac{0,0}{\infty}$} &      &      &    & \\
 точка, ($-2,378$) & & & \totop{0,45}& \totop{0,45} & \totop{0,2} & \totop{0,245}  \\
\hline
\end{tabular}
%\end{table}
%\end{minipage}
\end{sidewaystable}
В качестве матрицы планирования выбрана матрица центрального композиционного
ротатабельного униформ-планирования [5--7]. Планы такого рода для описания
областей, близких к оптимуму (почти стационарной области), предпочтительнее,
чем ортогональные планы, так как обладают свойством ротатабельности "---
равномерного распределения информации о процессе при движении от центра
плана в любом направлении. Отказ от ортогонального планирования связан с
определённым усложнением расчётов.

Использована ниже следующая структура плана.
\begin{itemize}
\item[1.] Полуреплика полного факторного эксперимента "--- $2^{k-1}$ ($k = 6$) [5--8], заданная генерирующим
соотношением $x_6 =x_1 \cdot x_2 \cdot x_3 \cdot x_4 \cdot x_5 $ и, следовательно, определяющим контрастом $1=x_1 \cdot x_2 \cdot x_3 \cdot x_4 \cdot x_5 \cdot x_6 $, содержащая 32 опыта ($x_i$ "--- кодовые обозначения
варьируемых факторов согласно табл.~2). Возможность использования этой полуреплики обусловлена тем, что в ней остаются несмешанными линейные эффекты и парные взаимодействия [5--7]. 
\item[2.] Ядро плана: двенадцать <<звёздных>> точек с длиной звездного плеча $\alpha  = 2,378$. 
\item[3.] Центр плана: девять точек в центре плана (ввиду детерминированности результатов вычислительного эксперимента реализована одна точка).
\item[4.] Общее количество численных опытов в плане "--- 53.
\end{itemize}

Условия реализации численных опытов представлены в табл.~2, на основе
которой составлена матрица планирования вычислительного эксперимента. Ввиду
большой размерности матрицы планирования (содержит 53 строки и 42 столбца)
она вместе с результатами численного эксперимента (значениями всех 13
параметров оптимизации $y_1 \div y_{13}$ ) в данной работе не
приведена.


В качестве параметров оптимизации, обозначенных $y_1$, $y_2$,  $\ldots$, $y_{13}$,
приняты все наиболее важные и разнородные по составу и структуре образования характеристики процесса контактного теплообмена. Соответствующие обозначения и краткая характеристика параметров оптимизации $y_1$, $y_2$, $\ldots$, $y_{11} $
даются ниже:
\begin{itemize}
\item[$y_1$]  $=\hat {f}_{s0}^{(\infty )} (0)_{\min}^r$ "--- безразмерная температура на гравюре инструмента в начальный момент контактирования на квазиустановившейся стадии;
\item[$y_2$]  $=\hat {T}_{fs}^{(\infty )} (0,{\tt Fo}_s^c )$ "--- безразмерная температура на гравюре в момент завершения этапа контактирования на квазиустановившейся стадии; 
\item[$y_3$]   $ =\hat {T}_{fz}^{(\infty )} (0,{\tt Fo}_s^c )$ "--- безразмерная температура контактной поверхности
заготовки в момент завершения этапа контактирования на квазиустановившейся стадии; 
\item[$y_4$]  $ ={\tt Ki}_z^{(\infty )} (0)$ "--- безразмерное значение теплового потока (критерия Кирпичёва) в
начальный момент контактирования на квазиустановившейся стадии;
\item[$y_5$]   $ ={{\tt Ki}_z^{(\infty )}(0)}/{{\tt Ki}_z^{(\infty )} ({\tt Fo}_s^c )}$ "--- коэффициент изменения
теплового потока на этапе контактирования квазиустановившейся стадии процесса; 
\item[$y_6$]  $ =\hat {T}_{fz}^{(\infty )} (1,{\tt Fo}_s^c )$ "--- безразмерная температура центра заготовки в момент завершения этапа контактирования на квазиустановившейся стадии; 
\item[$y_7$]  $ ={\hat {T}_{fz}^{(\infty )} (1,{\tt Fo}_s^c )}/{\hat {T}_{fz}^{(\infty )} (0,{\tt Fo}_s^c )}$ "--- коэффициент изменения температуры заготовки по сечению в момент завершения этапа контактирования на квазиустановившейся стадии; 
\item[$y_8$]  $=\hat {T}_{fz}^{(\infty )} (0,{\tt Fo}_s )_{\min } $ "--- минимальная температура контактной поверхности заготовки на этапе контактирования квазиустановившейся стадии;
\item[$y_9$]    $=\hat {T}_{fs}^{(\infty )} (0,{\tt Fo}_s )_{\max } $ "--- максимальная температура на гравюре на этапе контактирования квазиустановившейся стадии; 
\item[$y_{10}$]  $ ={\hat {T}_{fs}^{(\infty )} (0,0)}/{\hat {T}_{fs}^{(\infty )} (1,0)}$ "---
коэффициент изменения температуры инструмента по сечению в начальный момент контактирования на квазиустановившейся стадии процесса; 
\item[$y_{11}$]  $ ={\tt Ki}_z^{(1)} (0)$ "--- безразмерное значение теплового потока (критерий Кирпичёва) в начальный момент первого цикла.
\end{itemize}

Программой выполнения расчётов  предусмотрено также определение интенсивности теплообмена между заготовкой и штампом в
интервале контактирования первого цикла: 
\begin{itemize}

\item[$y_{12}$]  $={\tt Ki}_z^{(1)} (0)/{\tt Ki}_z^{(1)} ({\tt Fo}_s^c )$ "--- коэффициент изменения теплового потока на этапе контактирования первого цикла; 

\item[$y_{13}$] $={\tt Ki}_z^{(1)} (0)/{\tt Ki}_z^{(\infty )} (0)$ "--- коэффициент изменения теплового потока в начальный момент первого цикла ($m = 1$) и на квазиустановившейся стадии ЦКТ (характеризует отклонение процесса от режима единичного цикла).
\end{itemize}

Температурный режим единичного цикла при ЦКТ характеризуется тем,
что в конце каждого последующего цикла температурное поле
совпадает с исходным температурным распределением (перед началом
цикла).

Вычисление коэффициентов $b_0$, $b_i$, $b_{ij}$, $b_{ii}$ уравнений
регрессионного типа (полиномы второй степени)
\begin{equation}
\label{stulin:eq9}
y=b_0 +\sum\limits_{i=1}^k b_i x_i +
\sum\limits_{\begin{subarray}{l}
 i=1 \\
 i<j \\
 \end{subarray}}^k b_{ij} x_i x_j + \sum\limits_{i=1}^k b_{ii} x_i^2,
\end{equation}
где $k$ "--- число факторов, производилось по специальной программе, составленной в среде {\tt Microsoft Excel}. Для
этого использовалась предложенная выше матрица планирования и~соответствующие формулы \cite{stulin:bib5}:
\begin{gather}
\notag
b_0 =\frac{A}{N}\biggl[2\lambda
^2(k+1)\sum\limits_{j=1}^N y_j -2\lambda c\sum\limits_{i=1}^k
\sum\limits_{j=1}^N x_{ij}^2 y_j  \biggr]; \quad b_i
=\frac{c}{N} \sum\limits_{j=1}^N x_{ij} y_j; \\
\label{stulin:eq10}
b_{il} =\frac{c^2}{\lambda N}\sum\limits_{j=1}^N x_{ij} x_{lj} y_j
\quad (l>i); \\
\label{stulin:eq11} b_{ii} =\frac{A}{N}\biggl\{ c^2\bigl[
(k+2)\lambda -k \bigr]\sum\limits_{j=1}^N x_{ij}^2 y_j
+c^2(1-\lambda )\sum\limits_{i=1}^k \sum\limits_{j=1}^N x_{ij}^2
y_j -2\lambda c\sum\limits_{j=1}^N y_j \biggr\};
\end{gather}
где $k=6$, $N=53$, $n_0 =9$ ($n_0$ "--- число опытов в центре плана),
\begin{gather}
\notag
\lambda =\frac{kN}{(k+2)(N-n_0 )}=0,9034;\quad
c=N\cdot \biggl(\sum\limits_{j=1}^N x_{ij}^2 \biggr)^{-1}
=1,2237;\\
\label{stulin:eq12}
A=\bigl({2\lambda \left[ (k+2)\lambda -k
\right]}\bigr)^{-1}=0,4509.
\end{gather}

По разработанной методике реализации численного эксперимента построено 13
регрессионных моделей полиномиального типа, из которых в качестве
иллюстрации ниже приведено одно уравнение для $y_1$:
\begin{equation}
\label{stulin:eq13}
\begin{aligned}
 y_1 &=0,07605973-0,02882247x_1 +0,04878578x_2 +0,01539632x_3 - \\
 &-0,020589x_4 +0,01560497x_5 -0,03950619x_6 -0,00812144x_1 x_2 - \\
&-0,0064455x_1 x_3 +0,013485x_1 x_4 -0,00432237x_1 x_5 +\\
&+0,00056437x_1 x_6 + 0,00164187x_2 x_3 -0,00768225x_2 x_4 + \\
&+0,0088865x_2 x_5 -0,01374812x_2 x_6 - 0,00629544x_3 x_4 + \\
&+0,00859894x_3 x_5 -0,00587744x_3 x_6 -0,00386506x_4 x_5 - \\
&-0,00040231x_4 x_6 +0,00552994x_5 x_6 +0,01804052x_1^2 +\\
&+0,00721053x_2^2 + 0,0099892x_3^2 +0,00700192x_4^2 + \\
&+0,00793489x_5^2 +0,02447375x_6^2 .
\end{aligned}
\end{equation}

Оценка адекватности представления интересующих зависимостей
полиномами \eqref{stulin:eq9} выполнялась по обычной схеме
посредством применения множественного коэффициента корреляции \cite{stulin:bib5},
численные значения которого находились в интервале $[0,95; 1)$, что
свидетельствует о достаточной адекватности регрессионных моделей
полиномиальногого типа в планируемой области изменения факторов.

Таким образом, на основе разработанной методики планирования
численного эксперимента построены в безразмерном виде
математические полиномиальные модели изменения наиболее важных
характеристик, определяющих тепловое состояние системы <<заготовка~-- пограничный слой~-- инструмент>> на квазиустановившейся
стадии ЦКТ. Полученные модели аппроксимационного типа дают
возможность в наглядной и компактной форме всесторонне
проанализировать тепловой режим контактного взаимодействия двух
плоских тел в безразмерном многофакторном пространстве с учётом
взаимного влияния и значимости каждого из выбранных факторов, а
также определить оптимальные значения параметров.

\begin{thebibliography}{99}

\RBibitem{stulin:bib1}
\by Евдокимов~М.\,А., Стулин~В.\,В.
\paper Аналитическое решение задач циклического контактного теплообмена для системы двух плоских тел
\jour Вестн. Сам. гос. техн. ун-та. Сер. Физ.-мат. науки
\yr 2008
\issue 1(16)
\pages 119--129
\mathnet{http://mi.mathnet.ru/vsgtu583}
\misctransl
\by Evdokimov~M.\,A., Stulin~V.\,V.
\paper Analytic solution to cyclic contact heat exchange problems for a~system of two flat bodies
\jour Vestn. Samar. Gos. Tekhn. Univ. Ser. Fiz.-Mat. Nauki
\yr 2008
\issue 1(16)
\pages 119--129


\RBibitem{stulin:bib2}
\by Стулин~В.\,В., Крупко~Е.\,А.
\paper Приближенная аналитическая оценка теплового режима системы двух плоских тел в произвольном цикле контактного теплообмена
\jour Вестн. Сам. гос. техн. ун-та. Сер. Мат.
\yr 2009
\issue 2(10)
\pages 65--70
\misctransl
\by Stulin~V.\,V., Krupko~E.\,A.
\paper Approximate analytical evaluation of the thermal regime for two flat bodies in any cycle of heat contact 
\jour Vestn. Samar. Gos. Tekhn. Univ. Ser. Mat.
\yr 2009
\issue 2(10)
\pages 65--70


\RBibitem{stulin:bib3}
\by Стулин~В.\,В.
\paper Восстановление температуры и мощности тепловых источников на границе тел в пространстве преобразований
Лапласа
\jour Вестн. Сам. гос. техн. ун-та. Сер. Мат.
\yr 2007
\issue 2(6)
\pages 94--103
\misctransl
\by Stulin~V.\,V.
\paper Renewal of the temperature and power of heat sources on the bodies boundary in the space of Laplace transforms
\jour Vestn. Samar. Gos. Tekhn. Univ. Ser. Mat.
\yr 2007
\issue 2(6)
\pages 94--103


\RBibitem{stulin:bib4}
\by Самарский~А.\,А., Вабищевич~П.\,Н.
\book Вычислительная теплопередача
\publaddr М.
\publ Едиториал УРСС
\yr 2003
\totalpages 784
\misctransl
\by Samarskiy~A.\,A., Vabishchevich~P.\,N.
\book Computational Heat Transfer
\publaddr Moscow
\publ Editorial URSS
\yr 2003
\totalpages 784


\RBibitem{stulin:bib5}
\by Налимов~В.\,В., Чернова~Н.\,А.
\book Статистические методы планирования экстремальных экспериментов
\publaddr М.
\publ Наука
\yr 1965
\totalpages 340
\misctransl
\by Nalimov~V.\,V., Chernova~N.\,A.
\book Statistical Methods of Planning Extreme Experiments
\publaddr Moscow
\publ Nauka
\yr 1965
\totalpages 340

\RBibitem{stulin:bib6}
\by Адлер~Ю.\,П., Маркова~Е.\,В., Грановский~Ю.\,В.
\book Планирование эксперимента при поиске оптимальных условий
\publaddr М.
\publ Наука
\yr 1976
\totalpages 280
\misctransl
\by Adler Yu.\,P., Markova~E.\,V., Granovsky Yu.\,V.
\book The design of experiments to find optimal conditions
\publaddr Moscow
\publ Nauka
\yr 1976
\totalpages 280

\RBibitem{stulin:bib7}
\by Ахназарова~С.\,Л., Кафаров~В.\,В.
\book Методы оптимизации эксперимента в химической технологии
\publaddr М.
\publ Высш. шк.
\yr 1985
\totalpages 327
\misctransl
\by Akhnazarova~S.\,L., Kafarov~V.\,V.
\book Experiment Optimization in Chemistry and Chemical Engineering 
\publaddr Moscow
\publ Vyssh. shk.
\yr 1985
\totalpages 327

\RBibitem{stulin:bib8}
\by Кафаров~В.\,В.
\book Методы кибернетики в химии и химической технологии
\publaddr М.
\publ Химия
\yr 1976
\totalpages 464
\misctransl
\by Kafarov~V.\,V.
\book Cybernetic Methods in Chemistry and Chemical Technology
\publaddr Moscow
\publ Khimiya
\yr 1976
\totalpages 464

\Bibitem{stulin:bib9}
\by Hartmann K., Lezki E., Sch\"eafer W.
\book Statistische Versuchsplanung und -auswertung in der Stoffwirtschaft
\publaddr Leipzig
\publ VEB Deutscher Verlag f\"ur Grundstoffindustrie
\yr 1974
\totalpages 444
\rtransl русск. пер.:~
\by Хартман~К., Лецкий~Э., Шефер~В.
\book Планирование эксперимента в исследовании технологических процессов
\publaddr М.
\publ Мир
\yr 1977
\totalpages 552

\RBibitem{stulin:bib10}
\by Фихтенгольц~Г.\,М.
\book Курс дифференциального и интегрального исчисления
\publaddr М.
\publ Физматгиз
\yr 1962
\bookvol 2
\totalpages 808
\misctransl
\by Fikhtengolz~G.\,M.
\book A Course on the Differential and Integral Calculus
\publaddr Moscow
\publ Fizmatgiz
\yr 1962
\bookvol 2
\totalpages 808

\RBibitem{stulin:bib11}
\by Позняк~Л.\,А., Скрынченко~Ю.\,М., Тишаев~С.\,И.
\book Штамповые стали
\publaddr М.
\publ Металлургия
\yr 1980
\totalpages 244
\misctransl
\by Poznyak~A.\,L., Skrynchenko~Yu.\,M., Tishaev~S.\,I.
\book Die steels
\publaddr Moscow
\publ Metallurgiya
\yr 1980
\totalpages 244

\RBibitem{stulin:bib12}
\by Сторожев~М.\,В., Попов~Е.\,А.
\book Теория обработки металлов давлением
\publaddr М.
\publ Машиностроение
\yr 1977
\totalpages 423
\misctransl
\by Storoghev~M.\,V., Popov~E.\,A.
\book Theory of Metal Forming
\publaddr Moscow
\publ Mashinostroenie
\yr 1977
\totalpages 423


\RBibitem{stulin:bib13}
\by Тылкин~М.\,А., Васильев~Д.\,И., Рогалев~А.\,М., Шкатов~А.\,П., Бельский~Е.\,И.
\book Штампы для горячего деформирования металлов
\publaddr М.
\ed М.~А.~Тылкин
\publ Высш. шк.
\yr 1977
\totalpages 496
\misctransl
\by Tylkin~M.\,A. Vasiliev~D.\,I., Rogalev~A.\,M. Shkatov~A.\,P., Bel'skiy~E.\,I.
\book Stamps for hot deformation of metals
\publaddr Moscow
\ed M.~A.~Tylkin
\publ Vyssh. shk.
\yr 1977
\totalpages 496

\end{thebibliography}



\noindent
\hskip6.5cm \thedate \par\noindent
\hskip6.5cm\therevdate



%\newpage
%
%\chead[\fancyplain{}{\scriptsize \sl S\,t\,u\,l\,i\,n~V.\,V.}]
%{\fancyplain{}{\scriptsize \sl Regression Models Construction for Describing the Thermal System State \ldots}}

\makeengtitlem



 \newpage
%
% \tableofengcontents
%
%\end{document}
